% ===================================================================
%  TAULUKKO N:o 14 — Katu. — Kauppiaat.
%  Traduction 2026 d'apres texto/io, controlee sur texto/fr et
%  texto/sv. Consigne : voir texto/fi/10-taulukko-01.tex.
%
%  LE TABLEAU LE PLUS CHARGE DU LIVRET : quatre-vingt-dix-neuf
%  renvois, et le finnois y rencontre son genitif prepose deux fois.
%  « la aquo-spricilo (49) di la fonteno (48) » et « l'emperialo (59)
%  dil omnibuso (58) » : le possede porte un numero, le possesseur
%  aussi. Relative dans les deux cas.
%
%  AILLEURS LE GENITIF PREPOSE NE GENE PAS, parce que l'ido met deja
%  le possesseur en tete : « la magazino di la liquoristo (7) e ta di
%  la muzik-vendisto (10) » se traduit sans qu'un renvoi bouge.
% ===================================================================

%%K t14-apar-1 apar
\VUsaut{4.0mm}
\VUtitre{15.0pt}{60}{TAULUKKO N:o 14}
\VUsaut{3.0mm}

%%K t14-tit-1 sub
\VUpk{11.4pt}{40}{Katu. --- Kauppiaat.}
\VUsaut{3.0mm}

%%K t14-01-1 p
1. --- Ystäväni Juhani, joka asuu maalla, tuli viettämään kolme päivää perheeni
luona. \VUgras{Siihen} asti hänellä ei ollut ollut tilaisuutta nähdä suurta
kaupunkia; siksi käytämme kaikki päivämme kävelemiseen, ja minä selitän hänelle
sen, mitä hän ei tunne.

%%K t14-02-1 p
2. --- Ensin menimme katsomaan \VUgras{observatoriota} \textsuperscript{(1)},
joka kiinnosti häntä kovasti. Palatessamme sen \VUgras{kadun}
\textsuperscript{(3)} kautta, jonka varrella \VUgras{turkkilainen kylpylä}
\textsuperscript{(2)} on, kohtasimme \VUgras{hautajaissaattueen}
\textsuperscript{(4)}. \VUgras{Ruumissaaton} edellä kulki \VUgras{papisto},
joka edelsi \VUgras{ruumisvaunuja}, joihin \VUgras{arkku} oli asetettu. Monet
ystävät saattoivat \VUgras{surevaa} perhettä.

%%K t14-02-2 p
Saattueen kuljettua ohi ylitimme kadun suuren \VUgras{olutkapakan}
\textsuperscript{(5)} edessä, jossa monet \VUgras{asiakkaat} joivat
\VUgras{olutta} terassilla lasitetun \VUgras{markiisin} \textsuperscript{(6)}
suojaamina.

%%K t14-03-1 p
3. --- Juhani hämmästyi kovasti sitä \VUgras{kilpeä}, joka oli asetettu
\VUgras{liköörikauppiaan} \textsuperscript{(7)} puodin ja
\VUgras{musiikkikauppiaan} \textsuperscript{(10)} puodin väliin. Hän kysyi
minulta sen merkitystä; vastasin, että se osoittaa englantilaista
\VUgras{konsulaattia} \textsuperscript{(8)}, ja kiinnitin hänen huomionsa
\VUgras{konsuliin} \textsuperscript{(9)}, joka on parvekkeella.

%%K t14-tit-2 sub
\VUpk{10.2pt}{40}{Puotien tutkimista.}
\VUsaut{3.0mm}

%%K t14-04-1 p
4. --- Tietenkin pysähdyimme \VUgras{soittimien}, \VUgras{harmonien},
\VUgras{urkujen} ja \VUgras{pianojen} näyttelyn eteen; katselimme
\VUgras{nuottivihkojen} ja pianolle tarkoitettujen \VUgras{musiikkikappaleiden}
kuvitettuja kansia ja ihailimme kullatusta puusta tehtyä \VUgras{lyyraa}, jota
käytetään kylttinä.

%%K t14-05-1 p
5. --- Pölypilvet, joita \VUgras{kadunlakaisijat} \textsuperscript{(11)} ja
\VUgras{lakaisuvaunu} \textsuperscript{(12)} saivat aikaan, pakottivat meidät
kulkemaan nopeasti \VUgras{huonekalukauppiaan} \textsuperscript{(13)}
kauniiden \VUgras{huonekalustojen} ohi ja \VUgras{peltisepän}
\textsuperscript{(14)} \VUgras{uunien} ja \VUgras{lamppujen} näyttelyn ohi.

%%K t14-06-1 p
6. --- Muutoin meidän täytyi tuolla paikalla jättää jalkakäytävä, sillä se oli
täynnä paketteja, joita otettiin \VUgras{muuttoautosta} \textsuperscript{(16)}
pannakseen ne \VUgras{makasiiniin} \textsuperscript{(15)}, joka oli juuri
vuokrattu.

%%K t14-06-2 p
Se esti meitä näkemästä \VUgras{vaunusepän} \textsuperscript{(17)} kauniita
vaunuja, mutta ei estänyt meitä pysähtymästä katsomaan \VUgras{pakkaajaa}
\textsuperscript{(19)}, joka teki laatikkoa naapurilleen,
\VUgras{polkupyörien ja autojen kauppiaalle} \textsuperscript{(18)}. Sitten,
koska oli jo hieman myöhä, palasimme kotiin.

%%K t14-tit-3 sub
\VUpk{10.2pt}{40}{Kävely kaupungin läpi.}
\VUsaut{3.0mm}

%%K t14-07-1 p
7. --- Seuraavana päivänä äitini ehdotti Juhanille, että hänet valokuvattaisiin,
sitten hän antoi minulle tehtäväksi saattaa hänet \VUgras{valokuvaajan}
\textsuperscript{(20)} luo. Aamiaisen jälkeen menimme taiteilijan
\VUgras{ateljeehen}, jossa meidän täytyi odottaa jokseenkin kauan.
Askarruttaaksemme itseämme katselimme \VUgras{muotokuvia}
\textsuperscript{(22)}, jotka riippuvat seinällä.

%%K t14-07-2 p
Kun oli meidän vuoromme, työ ei kestänyt kauan, ja heti kun ystäväni oli
lopettanut poseeraamisen \VUgras{valokuvauskoneen} \textsuperscript{(21)}
edessä, laskeuduimme alas.

%%K t14-08-1 p
8. --- Ensimmäisessä kerroksessa näimme \VUgras{kampaajan}
\textsuperscript{(23)} ovesta, joka on auki, hänen kisällinsä, joka leikkasi
asiakkaan hiuksia.

%%K t14-08-2 p
Tultuamme kadulle kuljimme \VUgras{tupakkapuodin} \textsuperscript{(24)} ohi.
Juhania huvitti niin kovasti punainen \VUgras{lyhty} \textsuperscript{(25)} ja
\VUgras{paksu piippu}, jota käytetään \VUgras{kylttinä}
\textsuperscript{(26)}, että hän törmäsi kahteen vanhaan herraan, jotka
keskustelivat \VUgras{nuuskaa ottaen} \textsuperscript{(27)}.

%%K t14-tit-4 sub
\VUpk{10.2pt}{40}{Konditoria.}
\VUsaut{3.0mm}

%%K t14-09-1 p
9. --- Kaikkien maiden \VUgras{setelit} ja \VUgras{rahat}, jotka olivat
näytteillä \VUgras{rahanvaihtajan} \textsuperscript{(29)}
\VUgras{näyteikkunassa}, kiinnittivät muutaman hetken huomiotamme, mutta koska
oli välipalan aika, menimme \VUgras{konditorin} \textsuperscript{(31)} luo
viipymättä kauemmin.

%%K t14-10-1 p
10. --- Kun näimme kaikki \VUgras{herkut}, joita puoti sisälsi, \VUgras{kakut,
hunajaleivokset, korput} ja kaikenlaiset \VUgras{karamellit}, emme voineet
helposti valita. Lopulta Juhani valitsi \VUgras{kermakakut}, ja minä otin vain
pienen \VUgras{kirsikkatortun}, sillä makea \VUgras{koskee hampaisiini}, enkä
lainkaan halua käydä \VUgras{hammaslääkärin} \textsuperscript{(30)} luona
liian usein.

%%K t14-tit-5 sub
\VUpk{10.2pt}{40}{Konekirjoitus.}
\VUsaut{3.0mm}

%%K t14-11-1 p
11. --- Näyttääkseni ystävälleni \VUgras{kirjoituskoneen}, josta hän oli kuullut
puhuttavan mutta jota hän ei tuntenut, menimme \VUgras{konttoriin}
\textsuperscript{(32)}, joka kuuluu \VUgras{serkulleni} \textsuperscript{(34)}.
Tapasimme hänet sanelemassa kirjeitään \VUgras{sihteerilleen}
\textsuperscript{(35)}, joka on \VUgras{pikakirjoittaja}. Tuskin oli kirje
valmis, kun \VUgras{konekirjoittaja} \textsuperscript{(33)} jäljensi sen
koneellaan. Koska emme halunneet keskeyttää sukulaiseni toimia, viivyimme hänen
luonaan vain muutaman hetken.

%%K t14-12-1 p
12. --- Serkkuni konttorin alapuolella asuu suuri \VUgras{kelloseppä-jalokivikauppias}
\textsuperscript{(37)}, joka on lisäksi kultaseppä. Hänen komea puotinsa
sisältää monia \VUgras{kelloja, kaappikelloja, taskukelloja},
\VUgras{ketjuja} ja kallisarvoisia \VUgras{koruja}, esimerkiksi
\VUgras{helminauhoja}, kultaisia \VUgras{rannerenkaita},
\VUgras{korvarenkaita} ja \VUgras{sormuksia}, joita koristavat
\VUgras{jalokivet} ja \VUgras{timantit}. Yksi näyteikkunoista on erityisesti
varattu \VUgras{kultatöille} ja sisältää huomattavan kokoelman hopeisia
\VUgras{ruokailuvälineitä}.

%%K t14-13-1 p
13. --- Kääntyessämme kadunkulmasta huomasin, että se \VUgras{kilpi}
\textsuperscript{(36)} oli vaihdettu, joka on asetettu talon \VUgras{numeron}
viereen. Olin sanomaisillani huomioni ääneen, kun \VUgras{sen} sotilassoiton
iloiset sävelet kuuluivat. Aloimme juosta rykmenttiä vastaan ajattelematta, että
\VUgras{se} kulkisi \VUgras{hatuntekijän} \textsuperscript{(39)} ja
\VUgras{hansikkaantekijän} \textsuperscript{(40)} puotien ohi ja että olisimme
olleet yhtä hyvin sijoittuneet näkemään sen ohimarssin, jos olisimme jääneet
\VUgras{optikon} \textsuperscript{(38)} näyteikkunan eteen, joka on aivan täynnä
\VUgras{nenälaseja, lornetteja} ja \VUgras{kiikareita}.

%%K t14-tit-6 sub
\VUpk{10.2pt}{40}{Ohimarssi.}
\VUsaut{3.0mm}

%%K t14-14-1 p
14. --- \VUgras{Rykmentin} \textsuperscript{(41)} edellä kulki
\VUgras{rumpumajuri} \textsuperscript{(42)}, jota seurasivat
\VUgras{rummunlyöjät} \textsuperscript{(43)}, \VUgras{torvensoittajat}
\textsuperscript{(44)} ja koko \VUgras{soittokunta}. \VUgras{Kenraali}
\textsuperscript{(45)}, joka oli torilla adjutanttinsa kanssa, oli pysähtynyt
\VUgras{vesisuihkun} \textsuperscript{(49)} lähelle, joka nousee
\VUgras{suihkulähteestä} \textsuperscript{(48)}, ollakseen läsnä
\VUgras{ohimarssissa}.

%%K t14-14-2 p
\VUgras{Esikuntaupseerit} \textsuperscript{(46)}, \VUgras{eversti} ja
\VUgras{everstiluutnantti} tervehtivät häntä miekalla. \VUgras{Majurit} olivat
\VUgras{pataljooniensa} \textsuperscript{(47)} edessä ja kukin
\VUgras{kapteeni} komensi \VUgras{komppaniaansa}, samalla kun
\VUgras{luutnantit}, \VUgras{aliluutnantit} ja \VUgras{aliupseerit} pitivät
tavanomaista paikkaansa miestensä vieressä.

%%K t14-15-1 p
15. --- Monet ohikulkijat olivat kerääntyneet \VUgras{kadunristeykseen}
\textsuperscript{(50)}; kauppiaat, \VUgras{sateenvarjontekijä}
\textsuperscript{(51)}, \VUgras{värjäri} \textsuperscript{(53)},
\VUgras{aseseppä} \textsuperscript{(54)} tulivat ulos nähdäkseen
\VUgras{sotilaat}. Kaikki ajoneuvot, \VUgras{vuokravaunut}
\textsuperscript{(52)}, \VUgras{herrasvaunut} \textsuperscript{(57)} ja myös
\VUgras{kastelutynnyri} \textsuperscript{(56)}, asettuivat jalkakäytävän
viereen.

%%K t14-tit-7 sub
\VUpk{10.2pt}{40}{Kohtauksia kadulla.}
\VUsaut{3.0mm}

%%K t14-15-2 p
\VUgras{Kauppamatkustaja} \textsuperscript{(55)}, joka kantoi kädessään
\VUgras{näytelaukkujaan}, kulki nopeasti kadun yli, ja ne henkilöt, jotka
istuivat \VUgras{ylätasolla} \textsuperscript{(59)}, joka on
\VUgras{omnibussissa} \textsuperscript{(58)}, kääntyivät ympäri nauttiakseen
näytelmästä. Poloinen \VUgras{katulaulaja} \textsuperscript{(60)}
\VUgras{posetiiveineen} \textsuperscript{(61)} joutui keskeyttämään
\VUgras{laulunsa}, sillä rumpujen pauhu peitti hänen äänensä.

%%K t14-16-1 p
16. --- Kun rykmentti saapui \VUgras{kangaskaupan} \textsuperscript{(64)}
kohdalle, \VUgras{ompelijattaret} \textsuperscript{(63)} ja \VUgras{räätälit}
\textsuperscript{(62)} jättivät \VUgras{ompelukoneensa} ja työnsä katsoakseen
ulos ikkunasta. \VUgras{Kauppias} \textsuperscript{(65)}, joka oli juuri
pannut uusia \VUgras{pukuja} \textsuperscript{(66)} \VUgras{näyttelyynsä},
joutui panettamaan sisään \VUgras{kankaat} \textsuperscript{(67)} ja
\VUgras{liinatavarat}, jotka olivat asetettuina jalkakäytävälle
\VUgras{viemärinaukon} \textsuperscript{(68)} lähelle, peläten että väkijoukko
kaataisi ne, ja jopa vanha \VUgras{kulkukauppias} \textsuperscript{(69)},
jolta portinvartijatar tinki sukkia, joutui menemään käytävään päättääkseen
kauppansa.

%%K t14-16-2 p
Saatoimme rykmenttiä aina kasarmille asti, \VUgras{ja}, hyvin tyytyväisinä
päiväämme, palasimme päivällisen hetkellä.

%%K t14-tit-8 sub
\VUpk{10.2pt}{40}{Erilaisia elinkeinoja ja ammatteja.}
\VUsaut{3.0mm}

%%K t14-17-1 p
17. --- Seuraavana päivänä isäni ehdotti meille paikallisen lehden kirjapainossa
käyntiä. Otimme sen innostuneina vastaan.

%%K t14-17-2 p
\VUgras{Kirjanpainaja} on myös \VUgras{kirjakauppias} \textsuperscript{(70)}
\VUgras{(= kirjojenmyyjä)}, \VUgras{kustantaja}, \VUgras{paperikauppias} ja
\VUgras{kirjansitoja}. Hän on sijoittanut ensimmäiseen kerrokseen lehden
\VUgras{toimituksen} \textsuperscript{(71)} ja myös \VUgras{kaiverruspajan},
jossa työskentelevät \VUgras{litografit} \textsuperscript{(72)} ja
\VUgras{kuparinkaivertajat} \textsuperscript{(73)}, lopuksi
\VUgras{latomon}, jossa ovat \VUgras{painotyöläiset} \textsuperscript{(74)}.
Varsinainen \VUgras{painohuone} \textsuperscript{(75)},
\VUgras{painokoneet} ja eri koneet ovat pohjakerroksessa.

%%K t14-tit-9 sub
\VUpk{10.2pt}{40}{Juhani tekee hyvin monta kysymystä.}
\VUsaut{3.0mm}

%%K t14-18-1 p
18. --- Tullessamme ulos kirjapainosta Juhani pysähtyi hyvin hämmästyneenä
pienen lasitetun laatikon eteen, joka oli asetettu pylvääseen
\VUgras{lyhyttavarakaupan} \textsuperscript{(77)} kohdalle, ja kysyi, mihin sitä
käytetään. Isäni selitti hänelle, että se on \VUgras{palohälytin}
\textsuperscript{(76)}. Muutoinkin kaikki oli, kaupungissa, ystävälleni
ihmeellistä, aina yksinkertaisesta \VUgras{kivisuihkulähteestä}
\textsuperscript{(78)} ja kevyestä \VUgras{lähettikolmipyörästä}
\textsuperscript{(80)}, jota ajaa livreepukuinen \VUgras{lähettipoika}
\textsuperscript{(79)}, \VUgras{postivaunuihin} \textsuperscript{(81)} asti.

%%K t14-19-1 p
19. --- Samalla kun isäni muutamaksi hetkeksi jätti meidät puhuakseen
\VUgras{veronkantajan} \textsuperscript{(83)} kanssa, joka oli pysähtynyt
keskustelemaan \VUgras{asianajajan} \textsuperscript{(82)} ja \VUgras{notaarin}
\textsuperscript{(84)} kanssa, huvitimme itseämme seuraamalla silmillämme
liikennettä kadulla. Mitä erilaisimmat henkilöt kulkivat ohitsemme :
\VUgras{konditorin kisälli} \textsuperscript{(85)}, joka kantoi korissa komeaa
\VUgras{piirakkaa}, tuli \VUgras{vaatekauppiaan} \textsuperscript{(86)} takana;
\VUgras{lyhdynsytyttäjä} \textsuperscript{(87)} keskusteli \VUgras{kaasumiehen}
\textsuperscript{(88)} kanssa, \VUgras{nunna} \textsuperscript{(89)}, joka piti
orpotyttöä kädestä, oli palaamassa luostariinsa.

%%K t14-tit-10 sub
\VUpk{10.2pt}{40}{Kotimatkallamme.}
\VUsaut{3.0mm}

%%K t14-20-1 p
20. --- Sillä hetkellä, kun isäni jälleen liittyi meihin, Juhani nauroi
ääneen nähdessään kahden aivan noessa olevan \VUgras{nuohoojan}
\textsuperscript{(91)} likaiset kasvot ja omituisen repaleisen puvun.

%%K t14-21-1 p
21. --- Halusin ostaa kasvin äidilleni \VUgras{kukkienmyyjättäreltä}
\textsuperscript{(92)}, mutta isäni kiinnitti huomioni siihen, että se olisi
hyvin painava kantaa mukana ja että olisi parempi hankkia \VUgras{appelsiineja}.
Lähestyimme \VUgras{myyjätärtä} \textsuperscript{(94)}, ja kun
\VUgras{kotiopettajatar} \textsuperscript{(93)}, joka oli valitsemassa
kasvatilleen, oli päättänyt ostoksensa, annoimme palvella itseämme.

%%K t14-22-1 p
22. --- Palatessamme kotiin kohtasimme useita tuttavia : perheeni vanhan
ystävättären, joka nojasi käsivarttaan \VUgras{seuraneitiinsä}
\textsuperscript{(95)}, siltojen ja teiden \VUgras{insinöörin}
\textsuperscript{(96)}, ja yhden serkuistani \VUgras{lastenhoitajineen}
\textsuperscript{(98)}.

%%K t14-22-2 p
Sitten kohtasimme äitini \VUgras{modistin} \textsuperscript{(97)} ja kaksi
kaupungille vierasta \VUgras{munkkia} \textsuperscript{(99)}, jotka tulivat
luoksemme kysyäkseen isältäni tietä asemalle.

\VUsaut{6.0mm}
\VUfilet{22mm}
